\section{Sharp global error and the limits of endpoint tests}
\label{sec:sharp-global}

Projected modified fields and Duhamel propagation of the density defect
are established in recent work~\citep{obzherin2026density,obzherin2026finite}.
The distinction needed here is between an upper bound that sums the
magnitudes of local defects and the norm of their signed, transported
sum.  The latter gives an exact coefficient for the actual recursive
Euler method and a converse to the acceleration condition.

Write $F_{s,t}$ for the exact flow from $s$ to $t$.  For $s\le t$, define
the transported acceleration and its accumulated defect by
\begin{align}
 A_{s,t}&=(DF_{s,t}\,\acc_s)\circ F_{t,s},\nonumber\\
 B_t^\sigma&=\int_0^t\sigma(s)A_{s,t}\,\dd s.
 \label{eq:transported-global-defect}
\end{align}
The integral belongs to $L^2(p_t;\R^d)$.  It transports vectors by the
derivative of the flow, not merely by composition with the inverse flow.

\begin{theorem}[Exact global Wasserstein coefficient]
\label{thm:sharp-global-euler}
Assume the smoothness, growth, spatial Lipschitz, and moment hypotheses
of Theorem~\ref{thm:distribution_order}, without imposing the acceleration
condition.  Consider grids indexed by $n$, with $H_n\to0$, and scales $\eta_n\downarrow0$
such that $H_n=O(\eta_n)$.  On the $k$th interval set
$\sigma_n(s)=h_{n,k}/\eta_n$, and suppose
$\sigma_n\rightharpoonup^*\sigma$ in $L^\infty(0,T)$.
Let $\widehat p_t^n$ be the law of the piecewise linear Euler
interpolation.  Then
\begin{equation}
 \sup_{0\le t\le T}\left|
 \frac{W_2(\widehat p_t^n,p_t)}{\eta_n}
 -\frac12\|\Pi_{p_t}B_t^\sigma\|_{L^2(p_t)}\right|\longrightarrow0.
 \label{eq:sharp-global-uniform}
\end{equation}
In particular,
\begin{equation}
 \lim_n\frac{\max_kW_2(\widehat p_k^n,p_{t_{n,k}})}{\eta_n}
 =\frac12\max_{0\le t\le T}\|\Pi_{p_t}B_t^\sigma\|_{L^2(p_t)}.
 \label{eq:sharp-global-maximum}
\end{equation}
The coefficient is continuous in $t$.  No density, fourth moment, or
regularity of the projected field is assumed.
\end{theorem}

The proof first obtains the classical variational error expansion in
$L^2(p_0)$, uniformly in time.  It then proves a uniform Wasserstein
tangent expansion using the frozen-flow lemma, spatial truncation, and
the second-moment tail bound.  This establishes the coefficient without
assuming differentiability of a velocity-to-density map; complete
proofs are in Appendix~\ref{app:sharp-global}.

A useful specialization fixes a positive $C^1$ allocation density
$\rho$ with $\int_0^T\rho=1$ and chooses
$\int_0^{t_k}\rho=k/N$.  Then $\eta_N=1/N$ and
$\sigma_N\to1/\rho$ uniformly.  With
$g_s^T=\Pi_{p_T}A_{s,T}$, the terminal coefficient is
\begin{equation}
 \lim_{N\to\infty}N W_2(\widehat p_T^{N,\rho},p_T)
 =\frac12\left\|\int_0^T\frac{g_s^T}{\rho(s)}\,\dd s
             \right\|_{L^2(p_T)}.
 \label{eq:sharp-schedule-coefficient}
\end{equation}
Cancellation occurs inside this norm.  An allocation rule based on an
integral of nonnegative local magnitudes generally optimizes an upper
bound rather than this exact coefficient.

\begin{theorem}[Characterization of second-order distribution accuracy]
\label{thm:global-ap-characterization}
Under the same hypotheses, the following are equivalent for a fixed
field $u$ on $[0,T]$:
\begin{enumerate}
\item $\nabla\cdot(p_t\acc_t)=0$ for every $t$.
\item There is a constant $C$ such that every finite grid satisfies
      $\max_kW_2(\widehat p_k,p_{t_k})\le CH^2$.
\item Along one sequence of uniform grids, with $h=T/N\to0$,
      $\max_kW_2(\widehat p_k,p_{kh})=o(h)$.
\item For every positive smooth allocation density $\rho$,
      $W_2(\widehat p_T^{N,\rho},p_T)=o(N^{-1})$.
\end{enumerate}
In item 4 it suffices to use the densities in any fixed, arbitrarily
small $C^1$ neighborhood of the uniform density $1/T$.
\end{theorem}

Consequently, a fixed field has a sharp alternative on uniform grids:
the acceleration condition gives $O(h^2)$ error at all nodes; if the
condition fails, the maximum over nodes is $c_u h+o(h)$ for a strictly
positive $c_u$.  A first-order terminal coefficient that vanishes by
cancellation can always be made positive by an arbitrarily small smooth
perturbation of the allocation density.  These assertions concern a
fixed field as the grid is refined; they do not assert a common positive
constant over a changing family of learned fields.

\begin{corollary}[Odd-dimensional location--scale order dichotomy]
\label{cor:odd-global-order}
Consider a path $p_t=\Law(m(t)+g(t)Z)$ covered by
Theorem~\ref{thm:nonlinear-odd-rigidity}.  If either $m$ or $g$ is
nonaffine in time, every fixed generator satisfying the hypotheses
above has maximum Euler distribution error $c_uh+o(h)$ on uniform
grids, with $c_u>0$.  If both are affine, the unique
acceleration-preserving generator is Euler-exact, and every other
fixed generator in the stated regularity class has $\Theta(h)$
maximum distribution error.
\end{corollary}
This applies, in particular, to odd-dimensional expanding Gaussian
paths with a nonaffine scale, and to the location--scale paths of
finite nondegenerate Gaussian mixtures specified in
Theorem~\ref{thm:nonlinear-odd-rigidity}.

\begin{proposition}[Any finite set of terminal schedules can miss the defect]
\label{prop:finite-schedule-blindness}
Fix $a>0$ and any finite collection of positive smooth allocation
densities $\rho_1,\ldots,\rho_m$ on $[0,1]$.  There is a smooth field
generating the same path
$p_t=\mathcal N(0,e^{2at}I_2)$ as $a(I+J)x$ such that the acceleration
condition fails, yet
\[
 W_2(\widehat p_1^{N,\rho_i},p_1)=O(N^{-2})
 \qquad (i=1,\ldots,m).
\]
Its maximum distribution error on uniform grids is $\Theta(N^{-1})$.
The field can be arbitrarily close to $a(I+J)x$ in its time-dependent
matrix coefficients and any prescribed finite number of their
derivatives.
\end{proposition}

The construction is explicit up to a finite homogeneous linear system.
Choose a nonzero smooth $f$ with $\|f\|_\infty<a^2$ and
$\int_0^1 f/\rho_i=0$ for every $i$, and put
\begin{equation}
 u_t(x)=(aI+b(t)J)x,\qquad b(t)=\sqrt{a^2-f(t)}.
 \label{eq:blind-gaussian-generator}
\end{equation}
Here $J$ is planar rotation by $\pi/2$.  The projected acceleration is
exactly $f(t)x$.  Finitely many integral constraints leave nonzero
choices of $f$; the proof uses disjoint smooth bumps.  The same argument
can hide the first-order coefficient at any finite set of observation
times in addition to the terminal time.

\paragraph{Matching even the terminal second-order coefficient.}
Take $f(t)=\epsilon\cos(2\pi t)$ with $0<\epsilon<a^2$ in
Eq.~\eqref{eq:blind-gaussian-generator}.  Uniform Euler has
\begin{equation}
 \lim_{N\to\infty}N^2W_2(\widehat p_1^N,p_1)
 =\frac{2\sqrt2\,e^a a^3}{3},
 \label{eq:cosine-terminal-second}
\end{equation}
the same coefficient as the acceleration-preserving field $a(I+J)x$.
Nevertheless, along $N$ divisible by four,
\[
 \lim_N N W_2(\widehat p_{1/4}^N,p_{1/4})
 =\frac{\sqrt2\,e^{a/4}\epsilon}{4\pi}>0.
\]
For the arbitrarily small schedule perturbation
$\rho_\delta(t)=1+\delta\cos(2\pi t)$, $0<|\delta|<1$, the terminal
error instead satisfies
\begin{equation}
 \lim_N N W_2(\widehat p_1^{N,\rho_\delta},p_1)
 =\frac{\sqrt2\,e^a\epsilon}{2|\delta|}
       \left(\frac1{\sqrt{1-\delta^2}}-1\right)>0.
 \label{eq:cosine-perturbed-first}
\end{equation}
Thus even an accurate terminal convergence plot, including its leading
second-order coefficient, cannot by itself establish acceleration
preservation or second-order accuracy throughout the path.
